\documentclass[11pt,a4paper]{article}
\usepackage[T1]{fontenc}
\usepackage[utf8]{inputenc}
\usepackage{lmodern,amsmath,amssymb}
\usepackage[margin=25mm]{geometry}
\usepackage[hidelinks]{hyperref}
\hypersetup{pdftitle={A sharp mechanical cost for a finite-duration reversal},pdfauthor={navstokgap research working notes}}
\newcommand{\tightlist}{\setlength{\itemsep}{0pt}\setlength{\parskip}{0pt}}
\setlength{\emergencystretch}{3em}
\setcounter{secnumdepth}{0}
\begin{document}
\hypertarget{a-sharp-mechanical-cost-for-a-finite-duration-reversal}{%
\section{A sharp mechanical cost for a finite-duration
reversal}\label{a-sharp-mechanical-cost-for-a-finite-duration-reversal}}

A return with prescribed endpoint velocities \(+u\) and \(-u\) and
acceleration ceiling \(a\) requires duration at least \(2u/a\). Its
minimum kinetic action is \(mu^3/(3a)\), independent of any additional
waiting time. Meanwhile its sampled polygon action converges with a
sharp quadratic mesh bound. The endpoint reversal cost and the
refinement error therefore have different limits.

\hypertarget{model-and-sharp-minimum-c043}{%
\subsection{1. Model and sharp minimum
(C043)}\label{model-and-sharp-minimum-c043}}

Fix \(m,a,T>0\) and \(0<u<c\) in a specified inertial frame. The
admissible paths are \(X\in W^{2,\infty}(0,T)\), with continuous
velocity representative \(v=\dot X\),

\[
X(0)=X(T)=0,\quad v(0)=u,\quad v(T)=-u,\quad
|v|\le u,\quad |\dot v|\le a\quad\hbox{almost everywhere}.
\]

The functional is the kinetic cost \(S_K[X]=(m/2)\int_0^T v^2dt\), in
action units. The acceleration is an admissible external control, with
force \(m\dot v\). This is a bounded-control mechanics problem; a
prescribed autonomous potential, its potential contribution to
Hamilton's action and a dynamical momentum receiver would specify
additional physical data.

\textbf{Proposition.} The class is nonempty exactly when \(T\ge2u/a\).
For such \(T\), the unique minimizing velocity is, with \(\tau=u/a\),

\[
v_*(t)=\begin{cases}
u-at,&0\le t\le\tau,\\
0,&\tau\le t\le T-\tau,\\
-a(t-T+\tau),&T-\tau\le t\le T.
\end{cases}
\qquad S_{K,\min}=\frac{mu^3}{3a}.
\]

\textbf{Proof.} Lipschitz continuity gives \(2u=|v(T)-v(0)|\le aT\). For
a feasible duration the intervals \([0,\tau]\) and \([T-\tau,T]\) have
disjoint interiors. On the first, \(v(t)\ge u-at\ge0\); on the last,
\(v(t)\le-u+a(T-t)\le0\). Hence

\[
\int_0^T v^2dt\ge
\int_0^\tau(u-at)^2dt+
\int_{T-\tau}^T[u-a(T-t)]^2dt=\frac{2u^3}{3a}.
\]

The displayed \(v_*\) has zero total integral and therefore returns to
the starting position. It obeys both ceilings and attains equality.
Equality forces each endpoint ramp and zero velocity in the middle
almost everywhere; continuity gives uniqueness. Integration with
\(X(0)=0\) fixes the path. The maximum excursion is \(u^2/(2a)\),
reached at the waiting segment.

At \(T=2u/a\) the velocity is \(u-at\) throughout. Indeed equality in
the total velocity-change bound forces \(\dot v=-a\) almost everywhere,
so the whole admissible class is then a singleton. At longer durations
the admissible class includes variations even though the minimizing path
remains unique.

\hypertarget{which-premises-hold-the-cost-above-zero}{%
\subsection{2. Which premises hold the cost above
zero?}\label{which-premises-hold-the-cost-above-zero}}

The bound depends on prescribed nonzero endpoint speed and a finite
acceleration ceiling. At fixed force ceiling \(F_{\max}=ma\) it becomes

\[
S_{K,\min}=\frac{m^2u^3}{3F_{\max}}.
\]

For fixed \(m,a,T\), choosing successively smaller positive \(u\)
eventually preserves feasibility and gives \(S_{K,\min}\to0\). An upper
speed ceiling allows this family. Likewise, increasing the permitted
acceleration at fixed \(m,u,T\) drives the minimum to zero. Uniform
positivity over a preparation class requires corresponding lower bounds
on mass and endpoint speed and an upper bound on acceleration, or
another premise controlling their combination.

For \(T>2u/a\), choose a nonzero smooth function \(\phi\) supported
strictly inside the waiting interval, with \(\int\phi=0\). Such
functions are obtained by differentiating a nonconstant smooth compactly
supported bump. For small \(b\ne0\), \(v_b=v_*+b\phi\) satisfies the
same speed/acceleration bounds and all endpoint conditions. Since the
supports of \(\phi\) and the nonzero part of \(v_*\) are disjoint,

\[
S_K[v_b]-S_K[v_*]=\frac{mb^2}{2}\int\phi^2dt\longrightarrow0.
\]

Thus the positive minimum total cost coexists with positive excess costs
accumulating at zero. Its source is the endpoint-conditioned motion,
rather than a discrete spacing of admissible action values. The cost is
frame-specific: the endpoint velocities and return condition already
select that frame.

\hypertarget{the-sampled-polygon-error-c044}{%
\subsection{3. The sampled polygon error
(C044)}\label{the-sampled-polygon-error-c044}}

For any partition \(\pi:0=t_0<\cdots<t_N=T\), write \(d_i=t_{i+1}-t_i\),
\(|\pi|=\max_i d_i\), and define the chord kinetic action

\[
S_\pi=\frac m2\sum_i\frac{[X(t_{i+1})-X(t_i)]^2}{d_i}.
\]

For every admissible path, and more generally every velocity with
Lipschitz constant at most \(a\),

\[
0\le S_K-S_\pi\le\frac{ma^2}{24}\sum_i d_i^3
\le\frac{ma^2T}{24}|\pi|^2.
\]

\textbf{Proof.} The chord velocity on each cell is its average
\(\bar v_i\). Square completion and the pair-variance identity give

\[
S_K-S_\pi=\frac m2\sum_i\int_{t_i}^{t_{i+1}}(v-\bar v_i)^2dt,
\]

\[
\int_I(v-\bar v_I)^2dt
=\frac1{2d}\int_I\int_I[v(s)-v(t)]^2ds\,dt
\le\frac{a^2}{2d}\int_0^d\int_0^d(s-t)^2ds\,dt
=\frac{a^2d^3}{12}.
\]

Summing proves the claim. The constant \(1/24\) is sharp: the feasible
minimum-duration return has affine velocity with slope \(-a\) on every
cell, and attains the first upper bound for every partition. Equal-width
cells also attain the mesh bound. In the longer-duration minimizer,
cells contained inside a ramp saturate the local bound and cells inside
the waiting interval have zero error.

At fixed \(m,a,T\) this convergence is uniform over the admissible
class. For families with increasing acceleration bound \(a_N\), the
estimate still forces the defect to vanish if \(a_N|\pi_N|\to0\) (fixed
\(m,T\)). Maintaining a positive defect along shrinking meshes therefore
requires loss of this uniform acceleration control. This is a necessary
condition, not a construction of a positive remainder.

\hypertarget{source-comparison-and-next-mechanical-step}{%
\subsection{4. Source comparison and next mechanical
step}\label{source-comparison-and-next-mechanical-step}}

\href{https://liberzon.csl.illinois.edu/teaching/cvoc/node85.html}{Liberzon's
double-integrator example} provides the bounded-acceleration control
framework and a minimum-time problem with running cost one. Here the
endpoint data and running cost \(mv^2/2\) are different; the
endpoint-envelope proof above establishes the exact result.
\href{../references/batches/B21.md}{B21} records the bounded prior-art
search.

The next mechanical question is to realize finite-duration turns with a
specified conservative receiver or interaction, then derive its
correlation law. A07's external-control class provides a benchmark for
such a realization. In the companion G02 task, the receiver's internal
degrees of freedom also test which modes the available observables
detect.

Run \texttt{python3\ scripts/bounded\_acceleration\_checks.py} for exact
integrals, matching conditions, variations and rational partition tests.
The written proofs establish the general minima and uniform convergence.
\end{document}
